\documentclass[../../main.tex]{subfiles}

\begin{document}

\section{Radon-Nikodym 定理}\label{section:ch18-radon-nikodym}

令 $(X,\mathcal{M})$ 为可测空间. 对 $(X,\mathcal{M})$ 上的测度 $\mu$ 和 $X$ 上关于 $\mathcal{M}$ 可测的非负函数 $f$, 定义 $\mathcal{M}$ 上的集函数 $\nu$ 为对所有 $E\in\mathcal{M}$,
\begin{align}
\nu(E)=\int_{E}f\,d\mu.\label{align:ch18-rn-def}
\end{align}
我们从积分的线性和单调收敛定理推出, $\nu$ 是可测空间 $(X,\mathcal{M})$ 上的测度, 且它具有性质
\begin{align}
\text{若 }E\in\mathcal{M}\text{ 且 }\mu(E)=0,\text{ 则 }\nu(E)=0.\label{align:ch18-abs-cont-def}
\end{align}

\begin{definition}[绝对连续]\label{definition:ch18-abs-continuous}
测度 $\nu$ 称为关于 $\mu$ 绝对连续, 若\eqref{align:ch18-abs-cont-def}成立. 我们用记号 $\nu\ll\mu$ 表示 $\nu$ 关于 $\mu$ 绝对连续.
\end{definition}

\begin{proposition}\label{proposition:ch18-abs-cont-eps-delta}
令 $(X,\mathcal{M},\mu)$ 为测度空间, 而 $\nu$ 是可测空间 $(X,\mathcal{M})$ 上的有限测度. 则 $\nu$ 关于 $\mu$ 绝对连续当且仅当对每个 $\varepsilon>0$, 存在 $\delta>0$, 使得对任何集合 $E\in\mathcal{M}$,
\begin{align}
\text{若 }\mu(E)<\delta,\quad\text{则 }\nu(E)<\varepsilon.\label{align:ch18-eps-delta}
\end{align}
\end{proposition}

\begin{proof}
显然 $\varepsilon$-$\delta$ 准则\eqref{align:ch18-eps-delta}蕴涵 $\nu$ 关于 $\mu$ 绝对连续. 为证明逆命题, 我们用反证法. 假定 $\nu$ 关于 $\mu$ 绝对连续但 $\varepsilon$-$\delta$ 准则\eqref{align:ch18-eps-delta}不成立. 则存在 $\varepsilon_{0}$ 和 $\mathcal{M}$ 中的集列 $\{E_{n}\}$, 使得对每个 $n$, $\mu(E_{n})<1/2^{n}$ 且 $\nu(E_{n})\geqslant\varepsilon_{0}$. 对每个 $n$, 定义 $A_{n}=\bigcup_{k=n}^{\infty}E_{k}$. 则 $\{A_{n}\}$ 是 $\mathcal{M}$ 中集合的递减序列. 根据 $\nu$ 的单调性和 $\mu$ 的可数次可加性,
\begin{align*}
\text{对所有 }n,\quad\nu(A_{n})\geqslant\varepsilon_{0}\text{ 而 }\mu(A_{n})\leqslant 1/2^{n-1}.
\end{align*}
定义 $A_{\infty}=\bigcap_{n=1}^{\infty}A_{n}$. 根据测度 $\mu$ 的单调性, $\mu(A_{\infty})=0$. 由于 $\nu(A_{1})\leqslant\nu(X)<\infty$, 对所有 $n$, $\nu(A_{n})\geqslant\varepsilon_{0}$. 我们从测度 $\nu$ 的连续性推出 $\nu(A_{\infty})\geqslant\varepsilon_{0}$. 这与 $\nu$ 关于 $\mu$ 绝对连续矛盾.
\end{proof}

\begin{theorem}[Radon-Nikodym]\label{theorem:ch18-radon-nikodym}
令 $(X,\mathcal{M},\mu)$ 为 $\sigma$ 有限测度空间, 而 $\nu$ 是定义在可测空间 $(X,\mathcal{M})$ 上的 $\sigma$ 有限测度, 它关于 $\mu$ 绝对连续. 则存在 $X$ 上的非负函数 $f$, 它关于 $\mathcal{M}$ 可测, 使得对所有 $E\in\mathcal{M}$,
\begin{align}
\nu(E)=\int_{E}f\,d\mu.\label{align:ch18-rn-theorem}
\end{align}
函数 $f$ 在以下意义下是唯一的: 若 $g$ 是任何 $X$ 上也有该性质的非负可测函数, 则 $g=f$ a.e. $[\mu]$.
\end{theorem}

\begin{proof}
我们假设 $\mu$ 和 $\nu$ 都是有限测度而将推广到 $\sigma$ 有限的情形留作练习. 若对所有 $E\in\mathcal{M}$, $\nu(E)=0$, 则\eqref{align:ch18-rn-theorem}对 $X$ 上的 $f\equiv 0$ 成立. 因此假设 $\nu$ 在整个 $\mathcal{M}$ 上不消失.

我们首先证明存在 $X$ 上的非负可测函数 $f$ 使得
\begin{align}
\text{对所有 }E\in\mathcal{M},\quad\int_{E}f\,d\mu>0,\quad\int_{E}f\,d\mu\leqslant\nu(E).\label{align:ch18-rn-proof1}
\end{align}
对 $\lambda>0$, 考虑有限带号测度 $\nu-\lambda\mu$. 根据 Hahn 分解定理, 存在 $\nu-\lambda\mu$ 的 Hahn 分解 $\{P_{\lambda},N_{\lambda}\}$, 即 $X=P_{\lambda}\cup N_{\lambda}$ 且 $P_{\lambda}\cap N_{\lambda}=\varnothing$, 其中 $P_{\lambda}$ 是关于 $\nu-\lambda\mu$ 的正集而 $N_{\lambda}$ 是负集. 我们宣称: 存在某个 $\lambda>0$, 使得 $\mu(P_{\lambda})>0$. 假设不是这样. 令 $\lambda>0$, 则 $\mu(P_{\lambda})=0$. 因此 $\mu(E)=0$. 根据绝对连续性, $\nu(E)=0$. 由于 $N_{\lambda}$ 是关于 $\nu-\lambda\mu$ 的负集,
\begin{align}
\text{对所有 }E\in\mathcal{M}\text{ 与所有 }\lambda>0,\quad\nu(E)\leqslant\lambda\mu(E).\label{align:ch18-rn-proof2}
\end{align}
我们从这些不等式推出: 若 $\mu(E)>0$, $\nu(E)=0$. 当然, 根据绝对连续性, 若 $\mu(E)=0$, 则 $\nu(E)=0$. 由于 $\mu(X)<\infty$, 对所有 $E\in\mathcal{M}$, $\nu(E)=0$. 这是一个矛盾. 因此我们可以选取 $\lambda_{0}$ 使得 $\mu(P_{\lambda_{0}})>0$. 定义 $f$ 为 $\lambda_{0}$ 乘 $P_{\lambda_{0}}$ 的特征函数. 观察到 $\int_{X}f\,d\mu>0$ 且由于 $\nu-\lambda_{0}\mu$ 在 $P_{\lambda_{0}}$ 上是正的, 对所有 $E\in\mathcal{M}$,
\begin{align*}
\int_{E}f\,d\mu=\lambda_{0}\mu(P_{\lambda_{0}}\cap E)\leqslant\nu(P_{\lambda_{0}}\cap E)\leqslant\nu(E).
\end{align*}
因此\eqref{align:ch18-rn-proof1}对 $f$ 的这个选取成立.

定义 $\mathcal{F}$ 为 $X$ 上使得对所有 $E\in\mathcal{M}$,
\begin{align*}
\int_{E}f\,d\mu\leqslant\nu(E)
\end{align*}
的非负可测函数族. 接着定义
\begin{align}
M=\sup\left\{\int_{X}f\,d\mu\,\middle|\,f\in\mathcal{F}\right\}.\label{align:ch18-rn-M}
\end{align}
我们证明存在 $f\in\mathcal{F}$ 使得 $\int_{X}f\,d\mu=M$, 而\eqref{align:ch18-rn-theorem}对任何这样的 $f$ 成立. 若 $g$ 和 $h$ 属于 $\mathcal{F}$, 则 $\max\{g,h\}$ 也如此. 事实上, 对任何可测集 $E$, 将 $E$ 分解为 $E_{1}=\{x\in X\mid g(x)<h(x)\}$ 与 $E_{2}=\{x\in X\mid g(x)\geqslant h(x)\}$ 的不交并且观察到
\begin{align*}
\int_{E}\max\{g,h\}\,d\mu=\int_{E_{1}}h\,d\mu+\int_{E_{2}}g\,d\mu\leqslant\nu(E_{1})+\nu(E_{2})=\nu(E).
\end{align*}
选取 $\mathcal{F}$ 中的序列 $\{f_{n}\}$ 使得 $\lim_{n\to\infty}\int_{X}f_{n}\,d\mu=M$. 我们假设 $\{f_{n}\}$ 在 $X$ 上是逐点递增的, 否则用 $\max\{f_{1},\cdots,f_{n}\}$ 代替每个 $f_{n}$. 对每个 $x\in X$, 定义 $f(x)=\lim_{n\to\infty}f_{n}(x)$. 我们从单调收敛定理推出 $\int_{X}f\,d\mu=M$, 从而 $f$ 属于 $\mathcal{F}$. 对所有 $E\in\mathcal{M}$, 定义
\begin{align}
\eta(E)=\nu(E)-\int_{E}f\,d\mu.\label{align:ch18-rn-eta}
\end{align}
根据假设 $\nu(X)<\infty$, 因此 $\int_{X}f\,d\mu\leqslant\nu(X)<\infty$. 根据积分的可数可加性, $\eta$ 是一个带号测度. 它是一个测度, 由于 $f$ 属于 $\mathcal{F}$, 且它关于 $\mu$ 绝对连续. 我们宣称在 $\mathcal{M}$ 上 $\eta\equiv 0$, 因此\eqref{align:ch18-rn-theorem}对 $f$ 的这个选取成立. 事实上, 若不是这样, 我们用刚才的方法(其中 $\nu$ 被 $\eta$ 代替)得出, 存在非负可测函数 $\hat{f}$ 使得
\begin{align*}
\text{对所有 }E\in\mathcal{M},\quad\int_{E}\hat{f}\,d\mu>0\text{ 且 }\int_{E}\hat{f}\,d\mu\leqslant\eta(E)=\nu(E)-\int_{E}f\,d\mu.
\end{align*}
因此 $f+\hat{f}$ 属于 $\mathcal{F}$, $\int_{X}[f+\hat{f}]\,d\mu>\int_{X}f\,d\mu=M$ 与 $f$ 的选取矛盾.

剩下要证明唯一性. 若存在两个可积函数 $f_{1}$ 与 $f_{2}$ 使得\eqref{align:ch18-rn-theorem}成立, 则根据积分的线性性, 对所有 $E\in\mathcal{M}$,
\begin{align*}
\int_{E}[f_{1}-f_{2}]\,d\mu=0.
\end{align*}
因此在 $X$ 上 $f_{1}=f_{2}$ a.e. $[\mu]$.
\end{proof}

在习题 59 中我们概述了归功于 John von Neumann 的 Radon-Nikodym 定理的另一个证明: 它依赖于关于 Hilbert 空间的对偶的 Riesz-Fr\'echet 表示定理.

\begin{example}\label{example:ch18-rn-sigma-finite-necessary}
在 Radon-Nikodym 定理中 $\sigma$ 有限的假设是必要的. 事实上, 考虑可测空间 $(X,\mathcal{M})$, 其中 $X=[0,1]$ 而 $\mathcal{M}$ 是 $[0,1]$ 的 Lebesgue 可测子集族. 定义 $\mu$ 为 $\mathcal{M}$ 上的计数测度, 因此若 $E$ 是有限的, $\mu(E)$ 是 $E$ 中点的个数, 否则 $\mu(E)=\infty$. 仅有的测度为零的集是空集. 因此 $\mathcal{M}$ 上的每个测度关于 $\mu$ 绝对连续. 定义 $m$ 为 $\mathcal{M}$ 上的 Lebesgue 测度. 我们把证明不存在 $X$ 上的非负 Lebesgue 可测函数 $f$ 使得对所有 $E\in\mathcal{M}$,
\begin{align*}
m(E)=\int_{E}f\,d\mu
\end{align*}
留作练习.
\end{example}

回忆可测空间 $(X,\mathcal{M})$ 上的带号测度 $\nu$, 存在 Jordan 分解 $\nu=\nu_{1}-\nu_{2}$, 其中 $\nu_{1}$ 和 $\nu_{2}$ 是 $\mathcal{M}$ 上的测度, 它们中的一个是有限的. 我们定义测度 $|\nu|$ 为 $\nu_{1}+\nu_{2}$. 若 $\mu$ 是 $\mathcal{M}$ 上的测度, 带号测度 $\nu$ 称为关于 $\mu$ 绝对连续, 若 $|\nu|$ 关于 $\mu$ 绝对连续, 它等价于 $\nu_{1}$ 和 $\nu_{2}$ 都关于 $\mu$ 绝对连续. 从带号测度的这个分解以及 Radon-Nikodym 定理, 我们有以下关于有限带号测度的同样的定理.

\begin{corollary}\label{corollary:ch18-rn-signed}
令 $(X,\mathcal{M},\mu)$ 为 $\sigma$ 有限测度空间, 而 $\nu$ 是可测空间 $(X,\mathcal{M})$ 上的有限带号测度, 它关于 $\mu$ 绝对连续. 则存在函数 $f$, 它在 $X$ 上关于 $\mu$ 可积, 且对所有 $E\in\mathcal{M}$,
\begin{align*}
\nu(E)=\int_{E}f\,d\mu.
\end{align*}
\end{corollary}

回忆给定可测空间 $(X,\mathcal{M})$ 上的两个测度 $\mu$ 和 $\nu$, 我们说 $\mu$ 和 $\nu$ 是相互奇异的(写为 $\mu\perp\nu$), 若 $\mathcal{M}$ 中存在不交的集合 $A$ 和 $B$, 使得 $X=A\cup B$ 且 $\nu(A)=\mu(B)=0$.

\begin{theorem}[Lebesgue 分解定理]\label{theorem:ch18-lebesgue-decomposition}
令 $(X,\mathcal{M},\mu)$ 为 $\sigma$ 有限测度空间, 而 $\nu$ 是可测空间 $(X,\mathcal{M})$ 上的 $\sigma$ 有限测度. 则存在 $\mathcal{M}$ 上的测度 $\nu_{0}$(它关于 $\mu$ 奇异)以及 $\mathcal{M}$ 上的测度 $\nu_{1}$(它关于 $\mu$ 绝对连续), 使得 $\nu=\nu_{0}+\nu_{1}$. 测度 $\nu_{0}$ 和 $\nu_{1}$ 是唯一的.
\end{theorem}

\begin{proof}
定义 $\lambda=\mu+\nu$. 由于 $\mu$ 和 $\nu$ 是 $\sigma$ 有限测度, 测度 $\lambda$ 也是. 此外, $\mu$ 关于 $\lambda$ 绝对连续. Radon-Nikodym 定理告诉我们存在非负可测函数 $f$ 使得
\begin{align}
\text{对所有 }E\in\mathcal{M},\quad\mu(E)=\int_{E}f\,d\lambda=\int_{E}f\,d\mu+\int_{E}f\,d\nu.\label{align:ch18-ld-proof1}
\end{align}
定义 $X_{+}=\{x\in X\mid f(x)>0\}$ 与 $X_{0}=\{x\in X\mid f(x)=0\}$. 由于 $f$ 是可测函数, $X=X_{0}\cup X_{+}$ 是 $X$ 的可测集的不交分解, 因此 $\nu=\nu_{0}+\nu_{1}$ 是 $\nu$ 的作为相互奇异测度的和的表达式, 其中对所有 $E\in\mathcal{M}$,
\begin{align*}
\nu_{0}(E)=\nu(E\cap X_{0})\text{ 而 }\nu_{1}(E)=\nu(E\cap X_{+}).
\end{align*}
现在 $\mu(X_{0})=\int_{X_{0}}f\,d\lambda=0$, 由于在 $X_{0}$ 上 $f=0$, $\nu_{0}(X_{+})=\nu(X_{+}\cap X_{0})=\nu(\varnothing)=0$. 因此 $\mu$ 与 $\nu_{0}$ 相互奇异. 剩下来仅须证明 $\nu_{1}$ 关于 $\mu$ 绝对连续. 令 $\mu(E)=0$. 我们必须证明 $\nu_{1}(E)=0$. 然而, 由于 $\mu(E)=0$, $\int_{E}f\,d\mu=0$. 因此根据\eqref{align:ch18-ld-proof1}与积分在区域上的可加性,
\begin{align*}
\int_{E}f\,d\nu=\int_{E\cap X_{0}}f\,d\nu+\int_{E\cap X_{+}}f\,d\nu=0.
\end{align*}
但在 $E\cap X_{0}$ 上 $f=0$, 在 $E\cap X_{+}$ 上 $f>0$, 因此 $\nu(E\cap X_{+})=0$, 即 $\nu_{1}(E)=0$.
\end{proof}

使\eqref{align:ch18-rn-theorem}成立的函数 $f$ 称为 $\nu$ 关于 $\mu$ 的 Radon-Nikodym 导数. 它常记为 $\dfrac{d\nu}{d\mu}$.

\begin{exercise}\label{exercise:ch18-29}
证明关于有限测度 $\mu$ 和 $\nu$ 的 Radon-Nikodym 定理蕴涵关于 $\sigma$ 有限测度 $\mu$ 和 $\nu$ 的 Radon-Nikodym 定理.
\end{exercise}

\begin{exercise}\label{exercise:ch18-30}
证明 Radon-Nikodym 定理中函数 $f$ 的唯一性.
\end{exercise}

\begin{exercise}\label{exercise:ch18-31}
令 $\mu$, $\nu$ 以及 $\lambda$ 为可测空间 $(X,\mathcal{M})$ 上的 $\sigma$ 有限测度.
\begin{description}
\item[（i）] 若 $\nu\ll\mu$, 而 $f$ 是 $X$ 上关于 $\mathcal{M}$ 可测的非负函数, 证明
\begin{align*}
\int_{X}f\,d\nu=\int_{X}f\left[\frac{d\nu}{d\mu}\right]d\mu.
\end{align*}
\item[（ii）] 若 $\nu\ll\mu$ 且 $\lambda\ll\mu$, 证明
\begin{align*}
\frac{d(\nu+\lambda)}{d\mu}=\frac{d\nu}{d\mu}+\frac{d\lambda}{d\mu}\text{ a.e. }[\mu].
\end{align*}
\item[（iii）] 若 $\nu\ll\mu\ll\lambda$, 证明
\begin{align*}
\frac{d\nu}{d\lambda}=\frac{d\nu}{d\mu}\cdot\frac{d\mu}{d\lambda}\text{ a.e. }[\lambda].
\end{align*}
\item[（iv）] 若 $\nu\ll\mu$ 且 $\mu\ll\nu$, 证明
\begin{align*}
\frac{d\nu}{d\mu}\cdot\frac{d\mu}{d\nu}=1\text{ a.e. }[\mu].
\end{align*}
\end{description}
\end{exercise}

\begin{exercise}\label{exercise:ch18-32}
令 $\mu$, $\nu$, $\nu_{1}$ 以及 $\nu_{2}$ 为可测空间 $(X,\mathcal{M})$ 上的测度.
\begin{description}
\item[（i）] 证明: 若 $\nu\perp\mu$ 且 $\nu\ll\mu$, 则 $\nu=0$.
\item[（ii）] 证明: 若 $\nu_{1}$ 与 $\nu_{2}$ 关于 $\mu$ 奇异, 则对任何 $\alpha\geqslant 0$, $\beta\geqslant 0$, 测度 $\alpha\nu_{1}+\beta\nu_{2}$ 也如此.
\item[（iii）] 证明: 若 $\nu_{1}$ 与 $\nu_{2}$ 关于 $\mu$ 绝对连续, 则对任何 $\alpha\geqslant 0$, $\beta\geqslant 0$, 测度 $\alpha\nu_{1}+\beta\nu_{2}$ 也如此.
\item[（iv）] 证明 Lebesgue 分解中的唯一性断言.
\end{description}
\end{exercise}

\begin{exercise}\label{exercise:ch18-33}
证明: 在 Lebesgue 分解定理的证明中, 对所有 $E\in\mathcal{M}$,
\begin{align*}
\int_{E}g\,d\lambda=\int_{E}g\,d\mu+\int_{E}g\,d\nu.
\end{align*}
\end{exercise}

\end{document}